\documentclass[11pt, letterpaper]{article}
\usepackage[utf8]{inputenc}
\usepackage[T1]{fontenc}
\usepackage{lmodern}
\renewcommand{\familydefault}{\sfdefault}
\usepackage[spanish]{babel}
\usepackage[letterpaper, top=1.5cm, bottom=1.5cm, left=1.27cm, right=1.27cm]{geometry}
\usepackage{titlesec}
\usepackage{enumitem} 
\usepackage{xcolor}
\usepackage{booktabs}
\usepackage{tabularx}
\usepackage{array}
\usepackage{float}
\usepackage{amsmath}
\usepackage{longtable}
\usepackage{multirow}
\usepackage{graphicx}
\usepackage[hidelinks]{hyperref}

\linespread{1.1}

\titleformat{\section}{\Large\bfseries}{\thesection.}{0.5em}{}
\titleformat{\subsection}{\large\bfseries}{\thesubsection.}{0.5em}{}
\titleformat{\subsubsection}{\normalsize\bfseries}{\thesubsubsection.}{0.5em}{}
\titlespacing*{\section}{0pt}{18pt}{6pt}
\titlespacing*{\subsection}{0pt}{12pt}{4pt}
\titlespacing*{\subsubsection}{0pt}{8pt}{2pt}

\setlength{\parindent}{0pt}
\setlength{\parskip}{4pt}
\renewcommand{\arraystretch}{1.3}
\pagestyle{plain}

\begin{document}

\begin{titlepage}
    \centering
    \vspace*{1.5cm}

    {\Huge \textbf{SkyAnalytics Inc.}}\\[1cm]
    {\footnotesize\textbf{Drive:} \texttt{https://drive.google.com/drive/folders/1f9FfqWsVkM63CCfhutV755ouBE61dDS6?usp=sharing}}\\[0.2cm]
    {\footnotesize \textbf{Github:} \texttt{https://github.com/khrizzcoronel/SkyAnalytics.git}}\\[0.8cm]
    {\Large \textbf{Especificación de Requisitos de Software (SRS)}}\\[0.3cm]

    {\normalsize \textit{Basado en el Estándar IEEE 830}}\\[0.3cm]
    {\normalsize \textit{Khriz Coronel}}\\[0.3cm]
    {\normalsize \textit{Versión 1.0}}\\[0.3cm]

\end{titlepage}

\tableofcontents
\newpage

%=====================================================================
% 1. INTRODUCCIÓN
%=====================================================================

\section{Introducción}

\subsection{Propósito}
El propósito de este documento es definir la Especificación de Requisitos de Software (SRS) para la plataforma SaaS corporativa **SkyAnalytics**. Este documento sirve como contrato técnico detallado para guiar el desarrollo de la arquitectura híbrida (transaccional y analítica) del sistema, alineando las metas comerciales delineadas en el *Balanced Scorecard* con la ejecución técnica.

\subsection{Alcance}
El sistema **SkyAnalytics** es una plataforma unificada para la industria aeronáutica. Proporciona ingesta automatizada de datos de vuelo, modelos predictivos de Machine Learning, y tableros interactivos (*dashboards*) para el monitoreo estratégico y táctico. El sistema abarca:
\begin{itemize}
    \item \textbf{Capa Estratégica:} Visualización de KPIs de negocio, salud financiera y metas corporativas.
    \item \textbf{Capa Táctica (Command Center):} Monitoreo de integraciones de terceros (Sentry, GitHub), FinOps y reglas RBAC.
    \item \textbf{Capa Operativa (Headless):} Ingesta y limpieza masiva de datos (ETL), orquestación vía GitHub Actions, y despliegue efímero de entrenamiento ML (Instancias Spot).
\end{itemize}

\subsection{Definiciones, Acrónimos y Abreviaturas}
\begin{itemize}
    \item \textbf{SRS:} Software Requirements Specification.
    \item \textbf{BFF:} Backend-For-Frontend. Patrón arquitectónico donde el servidor actúa como capa intermedia dedicada a servir a un UI específico.
    \item \textbf{RBAC:} Role-Based Access Control.
    \item \textbf{ETL:} Extract, Transform, Load.
    \item \textbf{Zero Trust:} Filosofía de seguridad donde ninguna conexión de red es confiable por defecto (Ej. usar Túneles SSH para conexiones internas).
\end{itemize}

\subsection{Referencias}
\begin{itemize}
    \item Documento de Perfil Estratégico y Corporativo (\texttt{docs/analisis.tex}).
    \item Constitución de Arquitectura (SkyAnalytics \texttt{CONSTITUTION.md}).
    \item Especificaciones y Planes de Diseño Técnico (Módulos 001, 002 y 003).
\end{itemize}

%=====================================================================
% 2. DESCRIPCIÓN GENERAL
%=====================================================================

\section{Descripción General}

\subsection{Perspectiva del Producto}
SkyAnalytics se enmarca en un modelo de arquitectura \textbf{Monolito Modular Serverless}.
La aplicación principal se construye utilizando \textbf{Next.js (App Router)} y se despliega en un proveedor PaaS Serverless (Vercel). 
La persistencia de datos utiliza un enfoque híbrido:
\begin{itemize}
    \item \textbf{PocketBase (Operativa):} Base de datos ligera (SQLite) alojada en una VPS, responsable de la autenticación de usuarios, roles (RBAC), configuración del sistema, alertas y límites FinOps.
    \item \textbf{MonetDB (Analítica):} Motor columnar Data Warehouse especializado en procesamiento de grandes volúmenes de datos. Almacena las tablas de hechos (*Fact Tables*) alimentadas por el ETL y expone vistas materializadas para lectura instantánea.
\end{itemize}

\subsection{Funciones del Producto}
Las funciones del producto se dividen en 3 módulos principales:
\begin{enumerate}
    \item \textbf{Módulo Estratégico (Analítica Directiva):} Tableros de resumen (*Balanced Scorecard*), métricas financieras, telemetría y reportes de cumplimiento (exportables a PDF nativo).
    \item \textbf{Módulo Táctico (Command Center):} Gestión centralizada de accesos de usuario, rotación de secretos en el pipeline CI/CD, monitoreo de presupuesto Cloud (FinOps), y proxy de logs de Sentry.
    \item \textbf{Módulo Operativo (Backbone):} Trabajos programados vía GitHub Actions que ejecutan flujos ETL en Python (Pandas/Polars) con cuarentena de datos corruptos, y despliegue automatizado de modelos de ML en instancias efímeras (AWS Spot).
\end{enumerate}

\subsection{Características del Usuario}
El sistema está diseñado para 3 niveles de actores internos:
\begin{itemize}
    \item \textbf{\texttt{SUPER\_ADMIN}:} CTO o líder técnico. Tiene control total sobre el Módulo Táctico (secretos, presupuesto, roles).
    \item \textbf{\texttt{C\_LEVEL\_EXEC}:} Directivos y VP. Consumen el Módulo Estratégico para observar KPIs y establecer metas.
    \item \textbf{\texttt{OPERATOR}:} Ingenieros de Datos y ML. Diseñan y vigilan el Módulo Operativo.
\end{itemize}

\subsection{Restricciones y Suposiciones}
\begin{itemize}
    \item \textbf{Enfoque Solo-Dev / Bajo Presupuesto:} El sistema prohíbe el uso de infraestructuras pesadas permanentemente encendidas (ej. Apache Airflow o AWS RDS PostgreSQL multicapa).
    \item \textbf{Protección contra Timeouts:} Dado el entorno Serverless, las respuestas HTTP a APIs de terceros deben usar caché (\texttt{revalidate: 60}) y el entrenamiento de ML debe delegarse completamente a instancias externas.
\end{itemize}

%=====================================================================
% 3. REQUISITOS ESPECÍFICOS
%=====================================================================

\section{Requisitos Específicos}

\subsection{Requisitos de Interfaces Externas}
El sistema integrará y consumirá datos de las siguientes APIs externas como "Single Source of Truth":
\begin{itemize}
    \item \textbf{GitHub REST API:} Para lectura de estado de pipelines de CI/CD e inyección de Secrets de entorno.
    \item \textbf{Sentry API:} Para proxy en vivo de registros de errores.
    \item \textbf{Stripe API / Webhooks:} Para conciliación de pagos en tiempo real.
    \item \textbf{AWS Cost Explorer / PaaS Billing API:} Para auditoría financiera FinOps.
\end{itemize}

\subsection{Requisitos Funcionales por Módulo}

\subsubsection{Módulo Estratégico (Capa Directiva)}
\begin{table}[H]
\small
\centering
\begin{tabularx}{\linewidth}{>{\raggedright\arraybackslash}p{2.5cm} >{\raggedright\arraybackslash}X >{\raggedright\arraybackslash}p{2cm}}
\toprule
\textbf{ID Requisito} & \textbf{Descripción del Caso de Uso} & \textbf{Dependencia} \\
\midrule
RF-E01 (BSC) & El sistema debe renderizar el Balanced Scorecard con semaforización (Verde, Amarillo, Rojo) contra las metas LTM. & MonetDB \\
\midrule
RF-E02 (Finanzas) & El sistema debe calcular dinámicamente el ARR, Gross Margin, LTV, y CAC forzando la moneda USD. & MonetDB \\
\midrule
RF-E03 (Telemetría) & El sistema debe mostrar el SLA y Uptime de los servicios para auditar el \textit{Error Budget}. & MonetDB \\
\midrule
RF-E05 (Metas) & El \texttt{SUPER\_ADMIN} podrá configurar metas trimestrales vía CRUD. & PocketBase \\
\midrule
RF-E06 (Export) & Las vistas analíticas deben proveer un botón para invocar \texttt{window.print()} con CSS \texttt{@media print}. & N/A \\
\bottomrule
\end{tabularx}
\caption{Requisitos Funcionales del Módulo Estratégico}
\end{table}

\subsubsection{Módulo Táctico (Command Center)}
\begin{table}[H]
\small
\centering
\begin{tabularx}{\linewidth}{>{\raggedright\arraybackslash}p{2.5cm} >{\raggedright\arraybackslash}X >{\raggedright\arraybackslash}p{2cm}}
\toprule
\textbf{ID Requisito} & \textbf{Descripción del Caso de Uso} & \textbf{Dependencia} \\
\midrule
RF-T01 (RBAC) & Permitir definir matrices de permisos asociadas a Roles. & PocketBase \\
\midrule
RF-T02 (Logs) & Extraer y mostrar errores desde Sentry usando una tabla virtualizada (Infinite Scroll). & Sentry API \\
\midrule
RF-T03 (Secrets) & Proveer un formulario seguro para inyectar credenciales (Ej. DB\_PASSWORD) directo a GitHub Secrets. & GitHub API \\
\midrule
RF-T08 (FinOps) & Definir un presupuesto en PocketBase. Si el gasto del PaaS supera el 90\%, enviar Webhook a Slack. & AWS API \\
\bottomrule
\end{tabularx}
\caption{Requisitos Funcionales del Módulo Táctico}
\end{table}

\subsubsection{Módulo Operativo (Backbone)}
\begin{table}[H]
\small
\centering
\begin{tabularx}{\linewidth}{>{\raggedright\arraybackslash}p{2.5cm} >{\raggedright\arraybackslash}X >{\raggedright\arraybackslash}p{2cm}}
\toprule
\textbf{ID Requisito} & \textbf{Descripción del Caso de Uso} & \textbf{Dependencia} \\
\midrule
RF-O01 (Cron ETL) & GitHub Actions orquestará el pipeline ETL Python diario a las 02:00 AM. & GitHub Actions \\
\midrule
RF-O02 (Cuarentena) & Si el ETL encuentra datos corruptos, no debe crashear; debe desviarlos a la tabla \texttt{Quarantine\_Data}. & MonetDB \\
\midrule
RF-O03 (Backups) & El cron de backups debe generar un volcado cifrado y enviarlo a AWS S3. & S3/KMS \\
\midrule
RF-O04 (ML Spot) & El script ML levantará una instancia Spot, entrenará, subirá el modelo a S3 y ejecutará \texttt{shutdown -h now}. & AWS EC2 \\
\bottomrule
\end{tabularx}
\caption{Requisitos Funcionales del Módulo Operativo}
\end{table}

\subsection{Requisitos de Rendimiento}
\begin{itemize}
    \item \textbf{Tiempo de Respuesta (SLA):} Las llamadas al Backend-For-Frontend (BFF) del Módulo Estratégico deben resolverse en $\leq$ 3.0 segundos consultando las Vistas Materializadas de MonetDB.
    \item \textbf{Disponibilidad (Uptime):} La arquitectura tolerante a fallos debe apuntar a un Uptime del 99.99\%.
    \item \textbf{Caché en Proxy:} Las peticiones a APIs de terceros (Sentry, GitHub) en el Módulo Táctico deben ser cacheadas obligatoriamente durante al menos 60 segundos (\texttt{revalidate: 60}) para prevenir bloqueos por límite de cuota (Rate Limit 429).
\end{itemize}

\subsection{Atributos del Sistema de Software (Seguridad)}
Basados en las directivas de la Constitución (Zero Trust y Cifrado), el sistema impone las siguientes políticas:
\begin{itemize}
    \item \textbf{Zero Trust y Túneles:} Los servicios expuestos (ETL en GitHub Actions) no accederán al puerto abierto de MonetDB en internet público. Es obligatorio levantar un Túnel efímero (SSH / Tailscale) en el *Runner* antes de iniciar la ingesta.
    \item \textbf{Cifrado en Tránsito:} Toda comunicación hacia MonetDB requiere estrictamente \texttt{sslmode=require} (TLS 1.3).
    \item \textbf{Cifrado en Reposo:} El disco de la VPS que contiene el archivo SQLite de PocketBase y la Data de MonetDB debe estar cifrado a nivel de volumen (AWS EBS Encryption o LUKS AES-256). Los backups enviados a S3 deben usar KMS.
    \item \textbf{Autenticación Segura:} Las cuentas con roles administrativos requerirán soporte para MFA, interceptando accesos no autorizados en el *Middleware* de Next.js antes de la ejecución del código.
\end{itemize}

\end{document}
